% ============================================================================
% Fichier  : partie4/P04_C17_audit_production_continuite_tiers.tex
% Titre    : Auditer la Production, la Continuité et les Tiers
% Auteur   : MINKA MI NGUIDJOÏ Thierry Emmanuel
% Version  : 1.0 -- octobre 2026
% ============================================================================
% Sources : guide d'audit des SI du CHAI (2014), fiches 3.3 « production »
% (sauvegarde, restauration, livraisons, maintenance) et 3.5 « support et
% parc », paraphrasées ; COBIT 2019 (DSS01, DSS02, DSS04, BAI09) ; ISO
% 22301:2019 ; objectifs de reprise du système central de la BANQUE SANAGA
% (DMIA 8 h, RTO 4 h, RPO 15 min) et plan de continuité de 2022 jamais
% testé, repris du cours GCI (chapitre 15) ; panne du 9 mars 7 h 30 au
% 12 mars 8 h (72,5 h) repris du prologue de GCI. Rapport ISAE 3402 :
% chapitre 10. Données du journal de sauvegarde : base du ACSI-LAB-09.
% ============================================================================

\chapter{Auditer la Production, la Continuité et les Tiers}
\label{chap:audit-continuite}
\chaptermark{Production, continuité, tiers}

\epigraph{%
    « Une sauvegarde qu'on n'a jamais restaurée\\
    est une hypothèse, pas une sauvegarde. »%
}{Adage des exploitants}

\niveauChapitre{Intermédiaire}{%
    Maîtrise du Chapitre~\ref{chap:audit-securite} (infrastructure) et du
    Chapitre~\ref{chap:normes} (rapports ISAE~3402)%
}

\begin{boxobjectifs}
\begin{itemize}
    \item Décomposer l'audit de la production~: exploitation, sauvegardes et
          restaurations, support, parc, maintenance.
    \item Auditer des sauvegardes par leurs restaurations, et non par leurs
          journaux.
    \item Auditer un plan de continuité~: analyse d'impact, objectifs de
          reprise, tests, écart entre le visé et le constaté.
    \item Auditer un prestataire et un service en nuage~: clauses, droit
          d'audit, rapports d'assurance, dépendance.
\end{itemize}
\end{boxobjectifs}

\bigskip

% ============================================================================
\section*{Question directrice : la banque pouvait-elle reprendre en quatre heures~?}
% ============================================================================

Le plan de continuité de la BANQUE SANAGA, approuvé en 2022, fixe pour le
système central une durée maximale d'interruption admissible de 8~heures,
un objectif de reprise (RTO) de 4~heures et une perte de données maximale
(RPO) de 15~minutes (cours MSN-2S3-MSI03). Le 9~mars~2026, les agences
ouvrent sans système à 7~h~30~; le service revient le 12~mars à 8~h.
Personne n'a ouvert le plan.

La COBAC demande à Mme~EKAMBI si le dispositif de continuité était
capable, avant la panne, de tenir ses objectifs. Ce chapitre décompose
l'audit de la production (section~\ref{sec:c17-production}), audite les
sauvegardes (section~\ref{sec:c17-sauvegardes}) et la continuité
(section~\ref{sec:c17-continuite}), puis les prestataires
(section~\ref{sec:c17-tiers}). Le ACSI-LAB-09 applique la méthode aux journaux
de la banque.

\separateur

% ============================================================================
\section{Décomposer l'Audit de la Production}
\label{sec:c17-production}
% ============================================================================

La production, ou exploitation, fait fonctionner chaque jour les systèmes
que les projets ont livrés. Le guide du CHAI en propose deux fiches, l'une
pour la production, l'autre pour le support et le parc~\cite{chai2014guide}.

\begin{tableaucours}{cL{3.0cm}L{5.6cm}L{2.2cm}}{%
    \tableauHeader{Code} &
    \tableauHeader{Sous-domaine} &
    \tableauHeader{Contrôles (d'après le guide du CHAI)} &
    \tableauHeader{COBIT}
}
PX.1 & Exploitation & Traitements planifiés et suivis~; alertes remontées
en temps réel~; livraisons en production contrôlées par la production~;
guides d'exploitation à jour & DSS01 \\
PX.2 & Sauvegardes et restaurations & Sauvegardes conformes aux besoins~;
comptes rendus contrôlés~; supports externalisés~; restaurations testées
selon un planning & DSS04 \\
PX.3 & Support & Centre de services~; tickets qualifiés et escaladés~;
délais de clôture suivis~; satisfaction mesurée & DSS02 \\
PX.4 & Parc & Inventaire du matériel et des logiciels~; licences~;
installation à partir d'une configuration standard & BAI09 \\
PX.5 & Maintenance & Contrats de maintenance des matériels critiques et des
logiciels de base~; obsolescence suivie & DSS01, BAI09 \\
\end{tableaucours}
\vspace{-0.8em}
\begin{center}{\small\itshape Tableau~17.1. Décomposition de l'audit de la
production}\end{center}

Chez LiZbiZ, les domaines~B (sauvegardes) et~E (exploitation) du
Chapitre~\ref{chap:conclure} relèvent de cette décomposition~: le
domaine~E, en non-conformité opaque ($U = 0{,}30$), appelle d'abord la
recherche des preuves manquantes, notamment les comptes rendus des
traitements de nuit.

% ============================================================================
\section{Auditer les Sauvegardes}
\label{sec:c17-sauvegardes}
% ============================================================================

\subsection{Le journal ne suffit pas}

Un journal de sauvegarde dit qu'une copie a été écrite. Il ne dit pas
qu'elle est complète, lisible, ni qu'on saura la remonter dans le délai
voulu. Le guide du CHAI demande donc à l'auditeur de croiser les
sauvegardes avec les besoins, de vérifier le suivi des comptes rendus et
des supports, puis de « se focaliser » sur la restauration~: procédures
écrites, planning de tests, comptes rendus et actions qui en découlent.

\begin{tableaucours}{L{3.6cm}L{4.4cm}L{4.2cm}}{%
    \tableauHeader{Question} &
    \tableauHeader{Preuve} &
    \tableauHeader{Test}
}
Toutes les sauvegardes prévues ont-elles eu lieu~? & Journal de
sauvegarde, plan de sauvegarde & Rapprocher jour par jour le planifié et
l'exécuté \\
Les échecs sont-ils traités~? & Tickets, comptes rendus & Pour chaque échec,
retrouver l'action et sa date \\
Les supports sont-ils externalisés~? & Bordereaux de sortie & Rapprocher
les dates d'externalisation \\
Sait-on restaurer, et en combien de temps~? & Comptes rendus de
restauration & Comparer la durée constatée au RTO \\
\end{tableaucours}
\vspace{-0.8em}
\begin{center}{\small\itshape Tableau~17.2. Auditer les sauvegardes}\end{center}

\begin{boxalert}
\textbf{Violation typique.} Conclure que les sauvegardes sont efficaces
parce que le journal affiche « succès » chaque nuit.

\smallskip
\textbf{Règle~:} une sauvegarde est efficace quand une restauration
récente, documentée, a réussi dans le délai visé. Sans test de
restauration, le contrôle « on sait restaurer » est \NEK{}, quel que soit
le journal.
\end{boxalert}

% ============================================================================
\section{Auditer la Continuité}
\label{sec:c17-continuite}
% ============================================================================

\subsection{De l'analyse d'impact aux tests}

Le cours GCI présente la démarche de continuité~: analyse d'impact sur
l'activité, durée maximale d'interruption admissible (DMIA), objectifs de
reprise (RTO, RPO), plan, tests. L'ISO~22301 en formule les exigences pour
un système de management de la continuité~\cite{iso22301_2019}. L'auditeur
suit la même chaîne, et cherche à chaque maillon la pièce qui le prouve.

\begin{boxprocedure}
\textbf{Auditer un plan de continuité.}
\begin{enumerate}
    \item \textbf{Analyse d'impact}~: datée, revue après chaque changement
          majeur, approuvée par les métiers.
    \item \textbf{Objectifs}~: DMIA, RTO et RPO fixés pour chaque service
          critique, et cohérents (RTO inférieur à la DMIA).
    \item \textbf{Plan}~: il décrit le système \emph{actuel}, les
          personnes à joindre, les procédures de bascule et de retour.
    \item \textbf{Moyens}~: le site de secours héberge la version en
          production, et les données y sont répliquées au rythme du RPO.
    \item \textbf{Tests}~: au moins une fois par an pour chaque service
          critique, avec un compte rendu, la durée constatée et les actions.
    \item \textbf{Écart}~: la durée de reprise constatée, en test ou lors
          d'un incident réel, est comparée au RTO.
\end{enumerate}
\end{boxprocedure}

\subsection{Le dispositif de la BANQUE SANAGA}

\begin{boxlecture}{Continuité, système central}
Mme~EKAMBI cote dix contrôles de continuité du système central, à partir
du plan de 2022, des journaux de sauvegarde du premier trimestre~2026, des
comptes rendus de restauration de 2025 et de la chronologie de la panne.

\medskip
\begin{tableaucours}{L{7.0cm}c}{%
    \tableauHeader{Contrôle} &
    \tableauHeader{Cote}
}
Analyse d'impact revue depuis 2022 & \NON{} \\
RTO et RPO fixés et approuvés & \OUI{} \\
Plan à jour du système central actuel & \NON{} \\
Site de secours à la version de production & \NON{} \\
Plan testé dans l'année & \NON{} \\
RTO tenu lors du dernier incident & \NON{} \\
Sauvegardes quotidiennes réussies (seuil~: 98~\%) & \NON{} \\
Supports de sauvegarde externalisés & \OUI{} \\
Restauration testée chaque trimestre & \NON{} \\
Revue après la reprise formalisée & \NEK{} \\
\end{tableaucours}

\medskip
\textbf{Constatation.} $\mu = 0{,}20$~; $U = 0{,}10$~; $H \approx 1{,}16$~bit.
Le service a été indisponible 72,5~heures pour un RTO de 4~heures. Sur les
90~jours du premier trimestre, 84~sauvegardes ont réussi (93~\%). En 2025,
deux des quatre tests de restauration prévus ont eu lieu, dont un a
dépassé le RTO.

\medskip
\textbf{Analyse} (confiance élevée). Non-conformité claire~: la banque ne
pouvait pas tenir ses objectifs. Les objectifs existaient, les moyens ne
suivaient plus le système~: le plan et le site de secours décrivaient
l'ancien système central. Le seul \NEK{} porte sur la revue après
reprise~: le rapport à la COBAC existe, mais aucune revue interne
formalisée n'a été retrouvée.

\medskip
\hfill\lecture{0{,}20}{0{,}10}{1{,}16}
\end{boxlecture}

\begin{boxmetier}{Auditeur SI}
\textbf{Assister à un test de restauration}
\begin{enumerate}
    \item Choisir soi-même la sauvegarde à restaurer, et la date, dans la
          période auditée.
    \item Noter l'heure de début, les incidents, l'heure à laquelle le
          service est utilisable.
    \item Faire vérifier par un utilisateur métier que les données
          restaurées sont complètes (quelques opérations connues).
    \item Comparer la durée au RTO, et la date des dernières données
          restaurées au RPO.
    \item Conserver le compte rendu, signé, au dossier.
\end{enumerate}
\end{boxmetier}

% ============================================================================
\section{Auditer les Prestataires et le Nuage}
\label{sec:c17-tiers}
% ============================================================================

\subsection{On externalise une activité, pas sa responsabilité}

Infogérance, hébergement, service en nuage, maintenance applicative~: une
partie du SI est souvent exploitée par un tiers. Le risque reste à
l'organisation, et l'auditeur doit pouvoir l'apprécier. Trois sources de
preuve se complètent~: le contrat, les rapports d'assurance du prestataire
(un rapport ISAE~3402 de type~2, Chapitre~\ref{chap:normes}), et le droit
d'aller voir.

\begin{tableaucours}{L{4.0cm}L{8.2cm}}{%
    \tableauHeader{Clause} &
    \tableauHeader{Ce que l'auditeur vérifie}
}
Niveaux de service & Disponibilité, délais, RTO et RPO alignés sur ceux de
l'organisation~; pénalités appliquées \\
Droit d'audit & L'organisation, son auditeur et son régulateur peuvent
auditer le prestataire, ou obtiennent un rapport d'assurance équivalent \\
Sécurité et incidents & Mesures exigées~; notification des incidents dans
un délai fixé \\
Localisation des données & Lieu de stockage connu, compatible avec la loi
sur les données personnelles (Chapitre~\ref{chap:droit}) \\
Sous-traitance ultérieure & Recours à d'autres prestataires soumis à
accord \\
Réversibilité & Restitution des données dans un format exploitable, et
assistance à la sortie \\
Continuité du prestataire & Plan de continuité testé, dont les résultats
sont communiqués \\
\end{tableaucours}
\vspace{-0.8em}
\begin{center}{\small\itshape Tableau~17.3. Clauses à auditer dans un contrat
de prestation ou de service en nuage}\end{center}

\subsection{Le service en nuage}

Dans un service en nuage, les responsabilités sont \emph{partagées}~: le
fournisseur répond de l'infrastructure qu'il exploite, le client de ce
qu'il y configure (comptes, droits, données, paramètres de sécurité).
L'auditeur établit, pour chaque contrôle, qui en est responsable~; un
contrôle que chacun croit à la charge de l'autre n'est assuré par
personne. Les contrôles complémentaires énumérés dans un rapport ISAE~3402
jouent le même rôle~: ils restent à tester chez le client.

\begin{boxscenario}{Opération SANAGA. Le contrat de DIGICORE AFRIQUE}
\textbf{Faits.} Mme~EKAMBI lit le contrat de maintenance du système central
conclu avec DIGICORE AFRIQUE. Il prévoit des niveaux de service et des
pénalités, une clause de confidentialité et la localisation des données à
Douala. Il ne prévoit ni droit d'audit, ni réversibilité, ni encadrement de
la sous-traitance ultérieure. L'annexe sur la notification des incidents,
citée par le contrat, n'a pas été remise.

\medskip
\textbf{Cotation} des huit clauses~: 4~\OUI{}, 3~\NON{}, 1~\NEK{}.

\medskip
\textbf{Analyse} (confiance élevée). Sans droit d'audit, la banque dépend du
rapport que le prestataire choisit de produire~: un type~1, au
Chapitre~\ref{chap:normes}. Sans réversibilité, elle ne peut pas changer de
prestataire sans risque. La recommandation porte sur un avenant, à
négocier avant le renouvellement.
\end{boxscenario}

\begin{boxmetier}{Responsable GRC}
\textbf{Tenir le registre des prestataires critiques}
\begin{enumerate}
    \item Recenser les prestataires dont dépend un service critique, et
          le service concerné.
    \item Pour chacun, noter les clauses présentes, le dernier rapport
          d'assurance reçu, son type et sa période.
    \item Évaluer la dépendance~: existe-t-il une solution de remplacement,
          en combien de temps~?
    \item Porter au registre des risques tout prestataire critique sans
          droit d'audit ni réversibilité.
    \item Revoir le registre chaque année, et avant chaque renouvellement
          de contrat.
\end{enumerate}
\end{boxmetier}

% ============================================================================

\separateur

\begin{boxresume}
\textbf{Ce qu'il faut retenir de ce chapitre~:}
\begin{itemize}
    \item \textbf{Cinq sous-domaines de production}~: exploitation,
          sauvegardes et restaurations, support, parc, maintenance.
    \item \textbf{Une sauvegarde s'audite par sa restauration}~: sans test
          récent et documenté, « on sait restaurer » est \NEK{}.
    \item \textbf{La continuité est une chaîne}~: analyse d'impact,
          objectifs, plan, moyens, tests, écart entre le RTO et le
          constaté.
    \item \textbf{Un plan qui décrit l'ancien système} ne protège pas le
          nouveau~: il se met à jour à chaque changement majeur.
    \item \textbf{On externalise une activité, pas sa responsabilité}~:
          contrat, rapport d'assurance de type~2 et droit d'audit.
    \item \textbf{Dans le nuage}, chaque contrôle a un responsable
          désigné, fournisseur ou client.
\end{itemize}
\end{boxresume}

\bigskip

\begin{boxexo}[title={\ding{45}\quad Exercice~17.1. Rattacher les contrôles \niveaubase}]
Rattachez chaque contrôle à un sous-domaine (PX.1 à PX.5).
\begin{enumerate}[(a)]
    \item Les licences installées sont rapprochées des licences acquises.
    \item Les traitements de nuit ont un compte rendu contrôlé chaque matin.
    \item Le serveur de base de données est couvert par un contrat de
          maintenance avec intervention en quatre heures.
    \item Les tickets de priorité 1 sont escaladés en moins de 30~minutes.
    \item Un test de restauration est mené chaque trimestre.
\end{enumerate}
\end{boxexo}

\begin{boxexo}[title={\ding{45}\quad Exercice~17.2. RTO et RPO \niveaumoyen}]
Lors d'un test de restauration du système central de la BANQUE SANAGA, la
restauration commence à 9~h~00 et le service est utilisable à 14~h~10. La
dernière opération présente dans les données restaurées date de la veille
à 23~h~50~; la sauvegarde testée avait été lancée à 23~h~30.
\begin{enumerate}[(a)]
    \item Le RTO de 4~heures est-il tenu~? Calculez l'écart.
    \item Pourquoi ce test ne permet-il pas de juger le RPO de
          15~minutes~?
    \item Quel dispositif le RPO de 15~minutes suppose-t-il, et comment
          l'auditer~?
    \item Cotez le contrôle « RTO tenu en test » et rédigez la
          constatation.
\end{enumerate}
\end{boxexo}

\begin{boxexo}[title={\ding{45}\quad Exercice~17.3. Un service en nuage pour l'ANPP \niveauexpert}]
L'ANPP envisage d'héberger le portail de consultation des bulletins de
paie chez un fournisseur de service en nuage.
\begin{enumerate}[(a)]
    \item Dressez la liste des contrôles et, pour chacun, dites s'il
          relève du fournisseur ou de l'ANPP.
    \item Quelles clauses le contrat doit-il contenir~? Lesquelles sont
          liées au droit camerounais (Chapitre~\ref{chap:droit})~?
    \item Le fournisseur propose un rapport ISAE~3402 de type~1. Que
          demandez-vous, et pourquoi~?
    \item Comment la Chambre des comptes pourra-t-elle auditer ce service~?
\end{enumerate}
\end{boxexo}

\bigskip

\begin{boxinfo}
\textbf{Transition.} La Partie~\ref{part:pratique} a audité le SI domaine
par domaine. Le Chapitre~\ref{chap:cas-integrateurs} réunit les trois cas
du cours dans des missions complètes, de la lettre de mission au suivi.
\end{boxinfo}
